% Documento generado desde la guía DOCX equivalente.
% Compilación recomendada: pdflatex Guia_XX_*.tex (dos pasadas).
\documentclass[11pt,a4paper]{article}
\usepackage[utf8]{inputenc}
\usepackage[T1]{fontenc}
\usepackage[spanish,es-nodecimaldot]{babel}
\usepackage[a4paper,margin=2.35cm]{geometry}
\usepackage{lmodern}
\usepackage{microtype}
\usepackage[table]{xcolor}
\usepackage{array}
\usepackage{tabularx}
\usepackage{float}
\usepackage{enumitem}
\usepackage[most]{tcolorbox}
\usepackage{fancyhdr}
\usepackage{hyperref}
\definecolor{uemcgreen}{HTML}{004C3F}
\definecolor{uemcgold}{HTML}{E5B93F}
\definecolor{tfgblue}{HTML}{20566B}
\definecolor{tablelight}{HTML}{E7F1EE}
\definecolor{uemclight}{HTML}{F2F7F5}
\hypersetup{colorlinks=true,linkcolor=uemcgreen,urlcolor=tfgblue,pdftitle={Flujo y funcionamiento},pdfauthor={Alejandro Villarrubia García}}
\setlength{\parindent}{0pt}
\setlength{\parskip}{6pt}
\renewcommand{\arraystretch}{1.16}
\setlist[itemize]{leftmargin=1.2cm,itemsep=2pt,topsep=2pt}
\setlist[enumerate]{leftmargin=1.2cm,itemsep=4pt,topsep=2pt}
\pagestyle{fancy}
\fancyhf{}
\fancyhead[L]{\small\color{uemcgreen}GUÍA 01 · MAPA DEL SISTEMA}
\fancyhead[R]{\small\color{uemcgreen}TFG · UEMC}
\fancyfoot[C]{\color{uemcgreen}\thepage}
\renewcommand{\headrulewidth}{0.4pt}
\renewcommand{\headrule}{\hbox to\headwidth{\color{uemcgold}\leaders\hrule height \headrulewidth\hfill}}
\begin{document}
\begin{titlepage}
  \centering
  \vspace*{2.4cm}
  {\Large\color{uemcgreen}\textbf{GUÍA 01 · MAPA DEL SISTEMA}\par}
  \vspace{1.5cm}
  {\Huge\bfseries\color{uemcgreen}Flujo y funcionamiento\par}
  \vspace{0.7cm}
  {\large\color{tfgblue}Qué ocurre desde que entra un correo hasta que se explica el riesgo\par}
  \vfill
  {\large Proyecto TFG · versión alineada con el código actual\par}
  \vspace{0.35cm}
  {\large Alejandro Villarrubia García\par}
  \vspace{0.35cm}
  {\large Universidad Europea Miguel de Cervantes\par}
\end{titlepage}
\tableofcontents
\clearpage

\section{Idea general}

El montaje es cliente-servidor. El navegador usa Streamlit como capa de presentación y Streamlit envía peticiones HTTP/JSON al backend central. La extensión y el monitor también consumen esa API. Solo el servidor parsea, analiza, entrena y guarda los modelos; Gmail y Telegram son integraciones externas. La web comparte una interfaz adaptable con navegación de marca, estados homogéneos y un flujo de detección ordenado en tres pasos.

\begin{tcolorbox}[colback=uemclight,colframe=uemcgreen,title=Nota para la defensa]
\textbf{Frase para la defensa:} Aunque cliente y servidor estén en el mismo ordenador, son procesos separados con un contrato HTTP y comparten una versión activa por idioma. Para la demo móvil, mantén el backend en 127.0.0.1:8766 e inicia Streamlit con --server.address 0.0.0.0 --server.port 8501; el móvil usa la IP privada. Es acceso temporal a una LAN de confianza, no Internet, y la web no tiene autenticación multiusuario.
\end{tcolorbox}

\section{Flujo completo}

\begin{enumerate}
  \item Entrada de cliente: la web, la extensión o el monitor envían texto, campos JSON o un EML Base64 al backend; /analyze limita la petición a 16 MiB.
  \item Normalización: se extraen remitente, asunto, cuerpo, HTML, cabeceras completas, URLs, anclas y adjuntos.
  \item Señales: se revisan autenticación, coherencia de dominios, URLs, punycode, acortadores, HTML, formularios, lenguaje urgente y adjuntos.
  \item Puntuación heurística: RiskScorer pondera las señales activas y produce un riesgo de 0 a 100; ExplanationBuilder conserva el motivo.
  \item Idioma y modelo: EmailAnalysisService detecta español o inglés por mensaje y reutiliza un detector neuronal cacheado por idioma.
  \item Decisión: el modo combinado calcula una media ponderada; el umbral configurado se aplica de forma uniforme en heurístico, neuronal y combinado.
  \item Salida: el backend devuelve clasificación, puntuación, explicación, señales, idioma y versión; el cliente se limita a presentarlos o generar la alerta.
\end{enumerate}

\section{Piezas y responsabilidades}

\begin{table}[H]
  \centering
  \small
  \caption{Piezas y responsabilidades}
  \label{tab:comparativa-2}
  \begin{tabularx}{\textwidth}{|>{\raggedright\arraybackslash}p{2.45cm}|X|X|X|}
    \hline
    \rowcolor{uemcgreen}
    \textcolor{white}{\textbf{Capa}} & \textcolor{white}{\textbf{Módulo}} & \textcolor{white}{\textbf{Responsabilidad}} & \textcolor{white}{\textbf{Resultado}} \\
    \hline
    Entrada & analizador\_email.py & Parseo MIME seguro y extracción de partes & Correo normalizado \\
    \hline
    \rowcolor{tablelight}
    Dominio & correo.py & Contrato común para texto, diccionario o EML & CorreoAnalizado \\
    \hline
    Señales & header\_signals, url\_utils, html\_signals, content\_signals & Reglas pequeñas y comprobables & Diccionario de señales \\
    \hline
    \rowcolor{tablelight}
    Cliente & backend\_client.py & Contrato HTTP común sin lógica de detección & Petición/respuesta JSON \\
    \hline
    Servidor & backend\_service.py & Casos de uso, idioma, versiones y administración & Resultado central \\
    \hline
    \rowcolor{tablelight}
    ML & modelo\_neural.py & TF-IDF, MLP, entrenamiento y persistencia del servidor & Probabilidad phishing \\
    \hline
    Presentación & Streamlit, monitor, extensión & Recoger entrada y mostrar o alertar & UI o alerta \\
    \hline
  \end{tabularx}
\end{table}

\section{Cómo se analiza un EML}

El parser usa email.parser.BytesParser con una política MIME estándar. Conserva todas las cabeceras repetidas, no solo la última, y serializa las cabeceras originales junto al cuerpo. SPF, DKIM y DMARC se interpretan únicamente desde Received-SPF, Authentication-Results, ARC-Authentication-Results y DKIM-Signature: el prototipo no consulta DNS, no recupera claves públicas y no realiza una validación criptográfica completa.

\begin{itemize}
  \item Texto plano: se decodifica respetando el charset declarado y se conserva como cuerpo analizable.
  \item HTML: se guarda separado para revisar formularios, meta refresh, iframes, base href y javascript/data URLs.
  \item Adjuntos: se anotan nombre y tipo sin ejecutar ni abrir contenido potencialmente peligroso.
  \item URLs: se deduplican, se limpian de puntuación y se extraen también de enlaces HTML.
\end{itemize}

\section{Monitor y persistencia}

El monitor consulta Gmail, carga los IDs vistos y procesa solo mensajes nuevos. Cada correo normalizado se envía al backend con RemoteAnalysisService; el monitor no carga modelos. El estado se escribe después de cada mensaje completado mediante fichero temporal y os.replace.

Si un EML está corrupto o falla una alerta de Telegram, ese mensaje queda reportado con error y el resto del lote continúa. Un fallo de entrada no se convierte en una caída global del proceso.

\section{API y extensión}

\begin{table}[H]
  \centering
  \small
  \caption{Endpoints y controles de la API}
  \label{tab:comparativa-3}
  \begin{tabularx}{\textwidth}{|>{\raggedright\arraybackslash}p{2.45cm}|X|X|X|}
    \hline
    \rowcolor{uemcgreen}
    \textcolor{white}{\textbf{Endpoint}} & \textcolor{white}{\textbf{Entrada}} & \textcolor{white}{\textbf{Respuesta}} & \textcolor{white}{\textbf{Controles}} \\
    \hline
    GET /health y /models & Sin cuerpo & Estado y versiones sin rutas locales & No-store; sin artefactos \\
    \hline
    \rowcolor{tablelight}
    POST /analyze & JSON hasta 16 MiB & Riesgo, explicación, idioma y versión & Tipos, tamaño y CORS \\
    \hline
    POST /train y /evaluate & CSV serializado hasta 256 MiB & Versión y métricas & Token admin opcional en local \\
    \hline
    \rowcolor{tablelight}
    POST /compare y /models/delete & Configuraciones o idioma & Comparación o confirmación & Token y límites \\
    \hline
  \end{tabularx}
\end{table}

La extensión de Gmail no ejecuta Python: recoge el correo visible y lo envía directamente al backend central, por defecto en 127.0.0.1:8766. El proceso del puerto 8765 se conserva únicamente como proxy de compatibilidad y tampoco carga modelos.

\section{Ejemplo de extremo a extremo}

\begin{enumerate}
  \item El usuario abre un correo con asunto urgente y un enlace que aparenta ser de una marca conocida.
  \item El adaptador extrae From, Subject, cuerpo, anclas y URL.
  \item Las reglas detectan urgencia, dominio externo, anchor distinto y posible incoherencia de cabeceras.
  \item El modelo neuronal calcula una probabilidad sobre el texto completo.
  \item El modo combinado aplica la fusión calibrada 45/55 y el umbral 21; si una evidencia alcanza 70 no se diluye, y la explicación muestra qué señales activaron el riesgo.
  \item La UI pinta la tarjeta, la extensión la inserta en Gmail y el monitor podría notificar por Telegram.
\end{enumerate}

\section{Límites que conviene decir}

\begin{itemize}
  \item La lista de dominios y señales es local; no se afirma que exista reputación online en tiempo real.
  \item El análisis apoya decisiones y las 94 pruebas más 2 recorridos validan el flujo, no la eficacia en producción; el MLP requiere datos representativos y evaluación separada.
\end{itemize}

\end{document}
